\documentclass[11pt]{article}
\usepackage{coursenotes}
\begin{document}
\notesheader{1}{Causal basics}{Estimands, selection bias, randomization, and causal graphs}

\begin{decision}
A firm can send a coupon. It costs $c$ each time it is redeemed and each purchase earns margin $m$. The
number that settles the decision is the coupon's effect on purchase probability, compared with a
break-even that depends on who would have bought anyway. The question is how to get that effect out of data,
and what can go wrong on the way.
\end{decision}

\begin{goals}
  \item write a business question as an estimand in potential outcomes, and say which population it
  refers to;
  \item compute a treatment's net value and break-even, including the cost of ``sure things'';
  \item decompose a naive comparison into an effect plus selection bias, and say why randomization
  removes the bias;
  \item read a causal graph: chains, forks, colliders, and backdoor paths, and decide what to adjust for;
  \item show that the difference in means is unbiased over the randomization, and that it equals OLS on a
  dummy.
\end{goals}

%=====================================================================================
\sectionone

\subsection{The estimand}

Index units by $i = 1, \dots, n$ and let $D_i \in \{0,1\}$ record treatment. Each unit has two
\emph{potential outcomes}, $Y_i(1)$ if treated and $Y_i(0)$ if not; they exist before treatment is
assigned, and assignment only decides which one we see. The individual effect is
$\tau_i = Y_i(1) - Y_i(0)$.

\begin{defn}{Estimand}{estimand}
An estimand is a summary of potential outcomes for a named population. It specifies the units, the
intervention and its comparator, the outcome and its horizon. It is fixed before, and independently of,
any data or method.
\end{defn}

Three averages recur throughout the course.
\[
  \ATE = \E[Y(1) - Y(0)], \qquad \ATT = \E[Y(1) - Y(0) \mid D = 1], \qquad \ATU = \E[Y(1) - Y(0) \mid D = 0].
\]
The ATT and ATU depend on the assignment mechanism, because it defines who the treated are; the ATE does
not. Which one you need depends on the decision: ``keep treating the people we treat now?'' needs the ATT;
``treat everyone?'' needs the ATE; ``whom to treat?'' needs the conditional effect
$\tau(x) = \E[Y(1) - Y(0) \mid X = x]$, which later meetings learn to estimate.

\paragraph{Which population?} In the \emph{finite-population} (design-based) reading, $\E$ averages over
the $n$ units in the study and the estimand is the \emph{sample average treatment effect}
$\tau_S = \frac1n \sum_i \tau_i$, a fixed number. This is the reading these notes use for experiments,
because it is the answer about the customers in front of us and needs no assumption about how they were
sampled. In the \emph{superpopulation} reading, units are a sample from a larger population. The two
coincide when the study is a random sample of the population the decision covers; otherwise moving to
another population is a separate, assumed step called transport, which the course takes up later.

\begin{prop}{The ATE averages the ATT and ATU}{ate-mix}
With $\pi = \Pr(D=1)$, \quad $\ATE = \pi\,\ATT + (1-\pi)\,\ATU$.
\end{prop}
\begin{proof}
The law of total expectation applied to $Y(1) - Y(0)$, conditioning on $D$.
\end{proof}

\paragraph{Response types.} With a binary outcome each unit has one of four pairs $(Y(0), Y(1))$:
\emph{lost cause} $(0,0)$, \emph{sure thing} $(1,1)$, \emph{persuadable} $(0,1)$, and \emph{do-not-disturb}
$(1,0)$, with individual effects $0, 0, +1, -1$.

\begin{prop}{The ATE in response types}{types}
$\ATE = \Pr(\text{persuadable}) - \Pr(\text{do-not-disturb})$.
\end{prop}
\begin{proof}
$\E[\tau_i] = \sum_{\text{types}} \Pr(\text{type})\,\tau_{\text{type}}$, and only two types have a nonzero
effect.
\end{proof}

A lift of $0.10$ is consistent with 10\% persuadables and no do-not-disturbs, or with 30\% and 20\%. Data
identify the difference but not the two shares; Section~3 gives bounds.

\subsection{What a treatment is worth}

Let each purchase earn margin $m$, let the coupon cost $c$ per redemption (paid only when a treated unit
buys), and write $p_1 = \E[Y(1)]$, $p_0 = \E[Y(0)]$ for the population to be treated.

\begin{prop}{Net value and break-even}{econ}
The expected net value per unit treated, relative to not treating, is
\[
  V = m\tau - c\,p_1 = (m-c)\,\tau - c\,p_0,
\]
and treating pays if and only if $\tau > \tau^\ast = c\,p_1/m$.
\end{prop}
\begin{proof}
Treated, a unit yields $(m-c)Y(1)$; untreated, $m\,Y(0)$. The expected difference is
$(m-c)p_1 - m p_0 = m(p_1 - p_0) - c\,p_1$. Substitute $p_1 = p_0 + \tau$ for the second form.
\end{proof}

The second form separates what the coupon does: it earns $m-c$ on every purchase it causes and loses $c$
on every purchase that would have happened anyway, the \emph{sure-thing cost} $c\,p_0$. The table gives the
change per type, with $m = 30$ and $c = 10$.
\begin{center}
\begin{tabular}{lrrr}
\toprule
Type & Untreated & Treated & Change \\
\midrule
Persuadable & 0 & $m-c = 20$ & $+20$ \\
Sure thing & $m = 30$ & $m-c = 20$ & $-10$ \\
Lost cause & 0 & 0 & $0$ \\
Do-not-disturb & $m = 30$ & 0 & $-30$ \\
\bottomrule
\end{tabular}
\end{center}
A lost cause costs nothing when the treatment is paid only on use. If instead it costs $c$ per unit
\emph{treated} (a phone call, a trainer's hour), the net value is $m\tau - c$ and the break-even is
$\tau^\ast = c/m$. Always state which cost model applies. In both, the break-even depends on the
population, so a lift measured on one population must be held against the break-even of the population
the decision covers.

\subsection{Identification}

We observe $Y_i = D_i Y_i(1) + (1-D_i) Y_i(0)$, so one potential outcome is always missing: causal
inference is a missing-data problem in which the missing values are missing by design. Writing the
observed outcome this way already assumes two things, together called the stable unit treatment value
assumption (SUTVA).

\begin{assum}{SUTVA}{sutva}
(i) \emph{No interference}: unit $i$'s outcome depends only on its own treatment. (ii) \emph{Consistency}:
``treated'' names one well-defined intervention, so $Y_i = Y_i(d)$ when $D_i = d$.
\end{assum}

No interference fails when a coupon is shared with friends or a discount at one store draws shoppers from
the next. Consistency fails when ``a coupon'' is sometimes 10\% off and sometimes 30\%; then
the effect is a mixture over versions with weights set by whatever produced the data. In practice, describe
the intervention precisely enough that you could, in principle, assign it.

\begin{prop}{Selection-bias decomposition}{sb}
Under consistency,
\[
  \E[Y \mid D=1] - \E[Y \mid D=0] = \underbrace{\ATT}_{\text{effect on the treated}}
  + \underbrace{\E[Y(0) \mid D=1] - \E[Y(0) \mid D=0]}_{\text{selection bias}} .
\]
\end{prop}
\begin{proof}
By consistency the left side is $\E[Y(1) \mid D=1] - \E[Y(0) \mid D=0]$. Add and subtract
$\E[Y(0) \mid D=1]$.
\end{proof}

Selection bias asks how treated and untreated would have differed \emph{without} treatment. A manager who
treats units that would do well anyway makes it positive; one who treats struggling units makes it
negative and can reverse the sign. It does not shrink as the sample grows.

\begin{prop}{Two sources of bias relative to the ATE}{twobias}
Under consistency,
\[
\begin{aligned}
  \E[Y \mid D=1] - \E[Y \mid D=0] - \ATE &= \underbrace{\E[Y(0) \mid D=1] - \E[Y(0) \mid D=0]}_{\text{baseline selection}}\\
  &\quad + \underbrace{(1-\pi)(\ATT - \ATU)}_{\text{selection on gains}} .
\end{aligned}
\]
\end{prop}
\begin{proof}
Subtract $\ATE$ from Result~\ref{prop:sb} and use Result~\ref{prop:ate-mix}:
$\ATT - \ATE = (1-\pi)(\ATT - \ATU)$.
\end{proof}

\begin{assum}{Random assignment}{ra}
$D \indep \{Y(0), Y(1)\}$ and $0 < \Pr(D=1) < 1$.
\end{assum}

\begin{prop}{Identification by randomization}{id}
Under SUTVA and random assignment, $\ATE = \E[Y \mid D=1] - \E[Y \mid D=0]$, and $\ATE = \ATT = \ATU$.
\end{prop}
\begin{proof}
$\E[Y \mid D=d] = \E[Y(d) \mid D=d] = \E[Y(d)]$, by consistency and then independence. Independence also
makes both bias terms in Result~\ref{prop:twobias} zero.
\end{proof}

Independence is a property of the \emph{procedure}, not of the customers: a random number generator does
not know anyone's potential outcomes. That is why randomization removes confounding by variables nobody
recorded, and why it can be audited (the seed, the draw log, no post-draw exclusions) rather than argued
for. Independence holds over the distribution of possible draws. Any single draw can still be unbalanced by
chance, and standard errors measure how much that matters.

\subsection{Causal graphs}

A \emph{directed acyclic graph} (DAG) writes down assumptions about what causes what: an arrow $A \to B$
says $A$ may directly cause $B$; a missing arrow says it does not. The graph is a claim about the world,
not something estimated. Three building blocks decide whether two variables are associated.

\begin{center}
\begin{tikzpicture}[>={Stealth[length=2mm]}, every node/.style={font=\small}]
  \node (a1) at (0,0) {$A$}; \node (m1) at (1.3,0) {$M$}; \node (b1) at (2.6,0) {$B$};
  \draw[->] (a1) -- (m1); \draw[->] (m1) -- (b1); \node at (1.3,-1.0) {\textbf{Chain}};
  \node (a2) at (4.6,0) {$A$}; \node (f2) at (5.9,0.8) {$F$}; \node (b2) at (7.2,0) {$B$};
  \draw[->] (f2) -- (a2); \draw[->] (f2) -- (b2); \node at (5.9,-1.0) {\textbf{Fork}};
  \node (a3) at (9.2,0) {$A$}; \node (c3) at (10.5,-0.6) {$C$}; \node (b3) at (11.8,0) {$B$};
  \draw[->] (a3) -- (c3); \draw[->] (b3) -- (c3); \node at (10.5,-1.2) {\textbf{Collider}};
\end{tikzpicture}
\end{center}

\begin{itemize}
  \item \textbf{Chain} $A \to M \to B$, with $M$ a mediator: $A$ and $B$ are associated; conditioning on $M$
  blocks the association.
  \item \textbf{Fork} $A \leftarrow F \to B$, with $F$ a common cause: $A$ and $B$ are associated although
  neither causes the other; conditioning on $F$ blocks it.
  \item \textbf{Collider} $A \to C \leftarrow B$, with $C$ a common effect: $A$ and $B$ are \emph{not}
  associated; conditioning on $C$, or on anything it causes, \emph{creates} an association.
\end{itemize}

A \emph{path} is any sequence of arrows between two variables, whatever their direction. It is
\emph{blocked} by a conditioned-on middle node of a chain or fork, or by a collider that is not conditioned
on (nor any of its descendants). Two variables are associated, given a conditioning set, only through
unblocked paths.

\begin{prop}{Collider bias in the simplest case}{collider}
Let $D$ and $U$ be independent with $\Pr(D=1) = \Pr(U=1) = \tfrac12$, let $Y = U$ (so $D$ has no effect),
and let $C = \max(D, U)$. Then $\E[Y \mid D=1, C=1] - \E[Y \mid D=0, C=1] = \tfrac12 - 1 = -\tfrac12$.
\end{prop}
\begin{proof}
Given $D=1$, $C=1$ always, so the first mean is $\E[U] = \tfrac12$. Given $D=0$, $C=1$ forces $U=1$.
\end{proof}

A \emph{backdoor path} from $D$ to $Y$ is a path that starts with an arrow \emph{into} $D$.

\begin{prop}{Adjustment for a backdoor set}{backdoor}
Suppose a set $X$ of variables (i) contains nothing caused by $D$ and (ii) blocks every backdoor path from
$D$ to $Y$. Then conditional exchangeability holds, $\{Y(0), Y(1)\} \indep D \mid X$, and if
$0 < \Pr(D=1 \mid X) < 1$,
\[
  \E[Y(d)] = \sum_x \E[Y \mid D=d, X=x]\,\Pr(X=x) .
\]
\end{prop}
\begin{proof}
That (i) and (ii) imply conditional exchangeability is Pearl's backdoor theorem (Section~3). Given it,
$\E[Y(d)] = \sum_x \E[Y(d) \mid X=x]\Pr(X=x) = \sum_x \E[Y(d) \mid D=d, X=x]\Pr(X=x)$, and consistency
replaces $Y(d)$ by $Y$.
\end{proof}

The formula is \emph{standardization}: compare within levels of $X$, then average with the population's
weights; using the treated units' weights instead gives the ATT. Randomization is graph surgery. A coin
deletes every arrow into $D$, so no backdoor paths remain and no adjustment is needed.

\paragraph{Simpson's paradox.} An association can reverse when data are split by a third variable. Which
answer is right is decided by the graph: if the variable is a fork (a common cause of $D$ and $Y$), compare
within its levels and standardize; if it is a mediator or a collider, the pooled comparison is right and
splitting is wrong.

\subsection{Estimation: what the randomization guarantees}

In a completely randomized experiment exactly $n_1$ of $n$ units are treated, each subset equally likely.
Hold the potential outcomes fixed; the only randomness is the draw.

\begin{prop}{Unbiasedness over the design}{unbiased}
$\hatt = \bar Y_1 - \bar Y_0$ satisfies $\E_D[\hatt] = \tau_S$, where $\E_D$ averages over draws.
\end{prop}
\begin{proof}
$\bar Y_1 = \frac{1}{n_1}\sum_i D_i Y_i(1)$ and $\E_D[D_i] = n_1/n$ by symmetry, so
$\E_D[\bar Y_1] = \frac1n\sum_i Y_i(1)$; likewise $\E_D[\bar Y_0] = \frac1n\sum_i Y_i(0)$.
\end{proof}

The proof uses nothing but the draw: no outcome model, no sampling assumption, no large-sample
approximation. The collection of $\hatt$ values over all possible draws is the \emph{randomization
distribution}; its spread is what a standard error estimates.

\begin{prop}{OLS on a dummy is the difference in means}{ols}
In the least-squares fit of $Y_i = \alpha + \beta D_i + \varepsilon_i$, $\hat\alpha = \bar Y_0$ and
$\hat\beta = \bar Y_1 - \bar Y_0$.
\end{prop}
\begin{proof}
The normal equation for $\beta$ gives $\sum_{D_i=1} Y_i = n_1(\alpha + \beta)$, so $\alpha + \beta = \bar Y_1$;
the one for $\alpha$ then gives $\sum_{D_i=0} Y_i = n_0\alpha$.
\end{proof}

The regression is a computing device. Whether its coefficient is causal is decided by identification, not
by the software. The familiar omitted-variable-bias formula is Result~\ref{prop:sb} written for a regression: for binary
$D$, $\Cov(D,Y)/\Var(D) = \E[Y(0) \mid D=1] - \E[Y(0) \mid D=0] + \ATT$.

\subsection{Identification, estimation, inference}

Keep the three questions apart. \textbf{Identification}: with infinite data, would the comparison equal
the estimand? Settled by assumptions about how the data were generated. \textbf{Estimation}: how is the
comparison computed from finite data? \textbf{Inference}: how far is the estimate likely to be from the
truth because of chance? A sophisticated estimator cannot rescue a failed identification argument, and a
small standard error does not make a biased comparison right.

\begin{pitfall}[three common errors]
\begin{itemize}[leftmargin=1.2em]
  \item Holding a lift against zero instead of the break-even of the population the decision covers.
  \item ``Controlling for everything available'': a mediator removes part of the effect and a collider
  adds bias.
  \item Reading a large sample as protection against selection bias. It only makes the wrong number more
  precise.
\end{itemize}
\end{pitfall}

\checkpoint{A retailer finds that customers who used its app spent 40\% more than those who did not. Name
the potential outcomes, the estimand a decision to promote the app needs, the likely sign of selection bias,
and one fork in this setting.}

%=====================================================================================
\sectiontwo

\paragraph{Setting.} FitLife, a fitness-club chain, lets front-desk staff offer a free personal-training
session to members whose contracts are about to expire. The outcome $Y$ is renewal. A session costs ¥300
of trainer time whether or not the member renews, so the cost is per unit treated; a renewal is worth ¥1{,}000
of margin. Break-even: $\tau^\ast = 300/1{,}000 = 0.30$.

\paragraph{The hidden truth.} In a simulation both potential outcomes are visible. Staff offered sessions
to the four members they knew best, who mostly renew anyway.
\begin{center}
\begin{tabular}{lccccl}
\toprule
Member & $Y(0)$ & $Y(1)$ & $\tau_i$ & $D$ & Type \\
\midrule
A & 1 & 1 & 0 & 1 & sure thing \\
B & 1 & 1 & 0 & 1 & sure thing \\
C & 0 & 1 & 1 & 1 & persuadable \\
D & 1 & 1 & 0 & 1 & sure thing \\
E & 0 & 0 & 0 & 0 & lost cause \\
F & 0 & 1 & 1 & 0 & persuadable \\
G & 0 & 1 & 1 & 0 & persuadable \\
H & 1 & 1 & 0 & 0 & sure thing \\
\bottomrule
\end{tabular}
\end{center}

\paragraph{Estimands.} $\ATE = 3/8 = 0.375$; $\ATT = 1/4 = 0.25$; $\ATU = 2/4 = 0.50$. Check
Result~\ref{prop:ate-mix}: $0.5 \times 0.25 + 0.5 \times 0.50 = 0.375$.

\paragraph{The naive comparison.} Treated renewal rate $4/4 = 1.00$; untreated $1/4 = 0.25$; difference
$0.75$. Selection bias $= 3/4 - 1/4 = 0.50$, and $\ATT + 0.50 = 0.75$ (Result~\ref{prop:sb}). Relative to
the ATE: $0.50 + 0.5 \times (0.25 - 0.50) = 0.375$ of bias, so $0.375 + 0.375 = 0.75$
(Result~\ref{prop:twobias}).

\paragraph{The decision.} Taken at face value, the naive $0.75$ makes the program look like a clear success.
In truth the members staff choose gain $0.25 < 0.30$, so the program as run loses money; the members they
ignore gain $0.50$; and a blanket offer gains $0.375 > 0.30$. The selection that inflated the naive number also put the sessions in
the wrong hands. A randomized offer would have shown this.

\paragraph{A Simpson reversal at FitLife.} The chain also invited members to a free group class, mostly
new members, who renew less often.
\begin{center}
\begin{tabular}{lcccc}
\toprule
 & \multicolumn{2}{c}{Invited} & \multicolumn{2}{c}{Not invited} \\
\cmidrule(lr){2-3}\cmidrule(lr){4-5}
Member type & Renewed / members & Rate & Renewed / members & Rate \\
\midrule
New & 40 / 200 & 0.20 & 9 / 60 & 0.15 \\
Long-standing & 38 / 50 & 0.76 & 140 / 200 & 0.70 \\
\midrule
Pooled & 78 / 250 & 0.312 & 149 / 260 & 0.573 \\
\bottomrule
\end{tabular}
\end{center}
Pooled, invited members renew far less; within each type they renew more (by $0.05$ and $0.06$). Member
type is a fork. It drove invitations and renewal, and the class cannot change it. Standardizing to all 510
members (Result~\ref{prop:backdoor}) gives $(260 \times 0.05 + 250 \times 0.06)/510 = 0.055$; standardizing to
the invited members gives the ATT, $(200 \times 0.05 + 50 \times 0.06)/250 = 0.052$.

\begin{center}
\begin{tikzpicture}[>={Stealth[length=2mm]}, every node/.style={font=\small}]
  \node (t) at (0,1.2) {member type};
  \node (d) at (-2.2,0) {invited};
  \node (y) at (2.2,0) {renewed};
  \node (a) at (0,-1.1) {attended a class};
  \draw[->] (t) -- (d); \draw[->] (t) -- (y); \draw[->] (d) -- (y);
  \draw[->] (d) -- (a); \draw[->,dashed] (y) -- (a);
  \node[font=\footnotesize, text=extgrey] at (4.6,-1.1) {(dashed: members who plan to renew attend more)};
\end{tikzpicture}
\end{center}

\begin{exercise}[computation]
(a) Classify FitLife's eight members by response type and verify Result~\ref{prop:types}. (b) FitLife then
randomizes the session offer among 1{,}000 members and finds renewal rates $p_1 = 0.80$ and $p_0 = 0.45$.
Give the bounds on the persuadable share (Section~3), and say what extra assumption would pin it down.
\end{exercise}
\answer{(a) Persuadables C, F, G (3/8); sure things A, B, D, H; lost cause E; no do-not-disturbs. So
$\Pr(\text{persuadable}) - \Pr(\text{do-not-disturb}) = 3/8 = \ATE$. (b)
$\max(0, 0.80 - 0.45) = 0.35 \le \Pr(\text{persuadable}) \le \min(0.80, 0.55) = 0.55$. Monotonicity (no
member is put off by the offer) collapses the bounds to $0.35$.}

\begin{exercise}[judgment]
An analyst proposes to compare renewal among invited and uninvited members ``who attended at least one
class that month'', to compare like with like. Use the graph to explain what goes wrong.
\end{exercise}
\answer{Attendance is caused by the invitation and, through members' intention to renew, related to the
outcome: it is a descendant of $D$ and a collider on the path invited $\to$ attended $\leftarrow$ renewed.
Conditioning on it violates condition (i) of Result~\ref{prop:backdoor} and opens a non-causal path.
Uninvited members who attended anyway are unusually committed, so the comparison is biased toward zero or
below. Standardize on member type, which is pre-treatment, instead.}

%=====================================================================================
\sectionthree

\begin{extension}[Bounds on the persuadable share]
Randomization identifies $p_1$ and $p_0$, hence $\tau = \Pr(\text{P}) - \Pr(\text{DND})$, but not the
shares. Since $\Pr(\text{DND}) \ge 0$, $\Pr(\text{P}) \ge \tau$; since persuadables have $Y(1)=1$ and
$Y(0)=0$, $\Pr(\text{P}) \le \min(p_1, 1-p_0)$. So $\max(0, p_1 - p_0) \le \Pr(\text{P}) \le \min(p_1, 1-p_0)$.
The bounds collapse to $\tau$ under \emph{monotonicity} (no unit is harmed), an assumption about behavior
that the experiment cannot check. The same assumption returns later in the course as ``no defiers'', when
some customers do not take the treatment they are assigned.
\end{extension}

\begin{extension}[d-separation, the backdoor theorem, and beyond]
The blocking rules of Section~1.4 (causal graphs) are \emph{d-separation}. If a graph is correct, every d-separation it
implies appears as a conditional independence in the data, which gives a partial test (it can reject a
graph, never confirm one). Pearl's backdoor theorem shows that a set satisfying (i) and (ii) in
Result~\ref{prop:backdoor} delivers conditional exchangeability. When no such set is observed, the
\emph{front-door criterion} can identify an effect through a fully observed, unconfounded mediator.
Single-world intervention graphs (SWIGs) put potential outcomes on the graph and unify the two languages.
\end{extension}

\begin{extension}[Designs other than complete randomization]
Bernoulli randomization flips an independent coin per unit, so $n_1$ is random and the difference in means
is unbiased only conditional on both arms being non-empty. Stratified (blocked) designs randomize within
strata of a pre-treatment covariate; the estimator is the stratum-weighted average of within-stratum
differences, which is Result~\ref{prop:backdoor} applied where it is exact by construction.
\end{extension}

\begin{extension}[Attributes, audits, and causes of effects]
Attributes such as race or loyalty are not interventions; audit studies randomize a \emph{signal} of the
attribute (a name on comparable résumés) and estimate the effect of the signal on a reader's decision.
Reverse questions (``what caused this customer to leave?'') ask for causes of effects; they are best
rewritten as forward questions, one per candidate cause.
\end{extension}

\subsection*{Annotated reading}
\begin{readinglist}
\reading{Hern\'an, Miguel A., and James M. Robins. 2020. \emph{Causal Inference: What If}. Boca Raton: Chapman
\& Hall/CRC.}{Potential outcomes, randomization, standardization, and graphs in one consistent notation; free
online}{Ch.\ 1--3, 6--8}{1}
\reading{Pearl, Judea, Madelyn Glymour, and Nicholas P. Jewell. 2016. \emph{Causal Inference in Statistics: A
Primer}. Chichester, UK: Wiley.}{The shortest rigorous route to d-separation and the backdoor and front-door
criteria}{Ch.\ 1--3}{2}
\reading{Cinelli, Carlos, Andrew Forney, and Judea Pearl. 2024. ``A Crash Course in Good and Bad Controls.''
\emph{Sociological Methods \& Research} 53 (3): 1071--1104. \url{https://doi.org/10.1177/00491241221099552}}{A
catalog of graphs showing when adjusting helps or hurts}{All}{1}
\reading{Imbens, Guido W., and Donald B. Rubin. 2015. \emph{Causal Inference for Statistics, Social, and
Biomedical Sciences: An Introduction}. New York: Cambridge University Press.}{The design-based view of
experiments used in these notes}{Ch.\ 1--4}{2}
\reading{Holland, Paul W. 1986. ``Statistics and Causal Inference.'' \emph{Journal of the American Statistical
Association} 81 (396): 945--60. \url{https://doi.org/10.1080/01621459.1986.10478354}}{The classic statement of
the fundamental problem and ``no causation without manipulation''}{Sections 1--4}{1}
\reading{Bertrand, Marianne, and Sendhil Mullainathan. 2004. ``Are Emily and Greg More Employable Than Lakisha
and Jamal? A Field Experiment on Labor Market Discrimination.'' \emph{American Economic Review} 94 (4):
991--1013. \url{https://doi.org/10.1257/0002828042002561}}{How an audit study turns an attribute into a
randomized signal}{Intro, design}{1}
\end{readinglist}

\end{document}
